\documentclass[11pt,a4paper]{article}
\usepackage[utf8]{inputenc}
\usepackage[spanish]{babel}
\usepackage[margin=2.5cm]{geometry}
\usepackage{amsmath,amssymb,amsfonts}
\usepackage{graphicx}
\usepackage{booktabs}
\usepackage{listings}
\usepackage{xcolor}
\usepackage{fancyhdr}
\usepackage{ragged2e}
\usepackage{hyperref}

\definecolor{codegray}{rgb}{0.5,0.5,0.5}
\definecolor{backcolour}{rgb}{0.95,0.95,0.92}
\lstdefinestyle{mystyle}{
    backgroundcolor=\color{backcolour},   
    commentstyle=\color{codegray},
    keywordstyle=\color{blue},
    numberstyle=\tiny\color{codegray},
    stringstyle=\color{red},
    basicstyle=\ttfamily\footnotesize,
    breakatwhitespace=false,         
    breaklines=true,                 
    captionpos=b,                    
    keepspaces=true,                 
    numbers=left,                    
    numbersep=5pt,                  
    showspaces=false,                
    showstringspaces=false,
    showtabs=false,                  
    tabsize=2
}
\lstset{style=mystyle}

\pagestyle{fancy}
\fancyhf{}
\rhead{Proyecto 2: Motores de Inferencia Simbólica}
\lhead{Inteligencia Artificial -- Javeriana}
\rfoot{Página \thepage}

\begin{document}

\centering

% Baja la universidad un poco para que no quede pegada al encabezado
\vspace*{1.5cm}

% Bloque Superior: Institución
{\LARGE \textbf{Pontificia Universidad Javeriana}}\\[0.4cm]
{\large Facultad de Ingeniería -- Carrera de Ciencia de Datos}\\[0.3cm]
{\normalsize Introducción a la Inteligencia Artificial}

\vspace{0.8cm}

\vfill 

% Bloque Central: Título con interlineado controlado
{
    \huge \bfseries
    \begin{tabular}{c}
        Informe Técnico: Motor de Inferencia de \\
        Propósito General basado en Cláusulas de Horn
    \end{tabular}
}\\[0.8cm]
{\large \textit{Encadenamiento hacia Adelante y hacia Atrás en Dominios Médicos}}\\[1cm]
\textbf{Autores:}\\
María José Zabala Acero\\
Juan Andrés Bravo\\
David Alejandro Davila\\
Arturo Manrique\\[0.4cm]


\vfill 

% Bloque Inferior: Autores, Departamento y Fecha

\textit{Departamento de Ingeniería de Sistemas}\\[0.8cm]
Julio Omar Palacio Nino\\[2cm]
\textbf{29 de octubre de 2026}


\newpage



% --- A PARTIR DE AQUÍ EMPIEZA LA PÁGINA 2 ---
\justifying
\begin{abstract}
Este documento presenta el diseño, modelado e implementación de un motor de inferencia simbólica basado en lógica proposicional y cláusulas de Horn para un dominio médico complejo. Se detallan las estrategias algorítmicas de encadenamiento hacia adelante (\textit{forward chaining}) y hacia atrás (\textit{backward chaining}), garantizando una profundidad lógica de al menos cuatro niveles. Asimismo, se valida el funcionamiento del sistema mediante tres casos clínicos experimentales, demostrando la trazabilidad paso a paso en la deducción de diagnósticos y estados de alerta.
\end{abstract}

\vspace{0.5cm}

\section{Introducción y Definición del Problema}
El estudio de la inteligencia artificial clásica requiere la construcción de sistemas expertos capaces de resolver el problema de la consecuencia lógica $KB \models \alpha$, donde $KB$ representa la Base de Conocimiento estructurada en reglas de Horn de la forma $(P_1 \wedge P_2 \wedge \dots \wedge P_n \rightarrow Q)$, $\Delta$ denota el conjunto inicial de síntomas observados (hechos), y $\alpha$ es el diagnóstico u objetivo a verificar. 

Para cumplir con los requerimientos del proyecto, se ha desarrollado una base modular externa en formato JSON que contiene más de 50 reglas lógicas, estructuradas de lo general a lo particular a través de cuatro niveles de abstracción:
\begin{enumerate}
    \item \textbf{Nivel 1:} Hechos hacia Signos.
    \item \textbf{Nivel 2:} Signos hacia Síndromes.
    \item \textbf{Nivel 3:} Síndromes hacia Enfermedades.
    \item \textbf{Nivel 4:} Enfermedades y síntomas críticos hacia Subtipos, Tratamientos o Estados de Alerta.
\end{enumerate}
\section{Estrategias Algorítmicas Utilizadas}

\subsection{Motor de Inferencia hacia Adelante (Forward Chaining)}


\begin{itemize}
    \item \textbf{Funcionamiento:} 
    \item \textbf{Condición de Parada:}
    \item \textbf{Trazabilidad:} 
\end{itemize}

\subsection{Motor de Inferencia hacia Atrás (Backward Chaining)}


\begin{itemize}
    \item \textbf{Funcionamiento:} .
    \item \textbf{Casos Base y Control de Profundidad:} 
\end{itemize}

\section{Casos Clínicos de Prueba y Trazabilidad}

\subsection{Caso Clínico 1: }
\begin{itemize}
    \item \textbf{Síntomas Iniciales ($\Delta$):} \{\}
    \item \textbf{Diagnóstico Objetivo ($\alpha$):} N
    \item \textbf{Trazabilidad Paso a Paso (Hacia Adelante):}
    \begin{enumerate}
        \item \textit{Ciclo 1:} Se detectan Fiebre alta y Tos seca $\rightarrow$ Se infiere el signo intermedio \textit{Sind\_Respiratorio\_Leve}.
        \item \textit{Ciclo 2:} Con el síndrome anterior y Congestión nasal $\rightarrow$ Se infiere el síndrome \item \textit{Infeccion\_Vias\_Altas}.
        \item \textit{Ciclo 3:} Al incorporar Fatiga y persistencia de fiebre $\rightarrow$ Se deduce la enfermedad \textit{Neumonia\_Atipica} (Nivel 3).
        \item \textit{Ciclo 4:} Evaluando el diagnóstico con el síntoma crítico asociado $\rightarrow$ Se concluye el estado final \textit{Alerta\_Leve\_Monitoreo}. \textbf{¡Objetivo $\alpha$ alcanzado con éxito!}
    \end{enumerate}
\end{itemize}

\subsection{Caso Clínico 2:}
\begin{itemize}
    \item \textbf{Síntomas Iniciales ($\Delta$):} \{ hola \}
    \item \textbf{Diagnóstico Objetivo ($\alpha$):} 
    \item \textbf{Trazabilidad Paso a Paso:}
    \begin{enumerate}
        \item Se activa la rama general de cefaleas a partir de Cefalea intensa y Náuseas.
        \item El motor evalúa las bifurcaciones que comparten antecedentes similares y detecta el hecho diferenciador: \textit{Sensibilidad a la luz (Fotofobia)}.
        \item Gracias a este hecho diferenciador, el encadenamiento desvía la inferencia hacia la regla específica de \textit{Migraña con Aura}, descartando cuadros tensionales genéricos y alcanzando el nivel 3 de profundidad.
    \end{enumerate}
\end{itemize}

\subsection{Caso Clínico 3:}
\begin{itemize}
    \item \textbf{Síntomas Iniciales ($\Delta$):} \{\}
    \item \textbf{Diagnóstico Objetivo ($\alpha$):} 
    \item \textbf{Trazabilidad Paso a Paso (Hacia Atrás):}
    \begin{enumerate}
        \item El motor parte del objetivo $\alpha$: \textit{Codigo\_Rojo\_Urgencias}.
        \item Se evalúan los antecedentes requeridos: se necesita comprobar la existencia previa del síndrome \textit{Falla\_Cardiovascular} y el síntoma crítico \textit{Sudoración fría}.
        \item Al verificar hacia atrás, \textit{Falla\_Cardiovascular} se desglosa en \textit{Dolor torácico opresivo} y \textit{Dificultad respiratoria severa}, los cuales se encuentran en $\Delta$.
        \item Se valida la conjunción lógica, cerrando con éxito la prueba de consecuencia lógica por encadenamiento hacia atrás y generando la traza de auditoría clínica.
    \end{enumerate}
\end{itemize}

\section{Conclusiones}

\end{document}